\section{Method}
\label{sec:method}

\subsection{Same-time candidates from different observations}

Let $E_{u,q}\in\mathbb{R}^{7}$ denote the action for physical controller time
$u$ predicted by the frozen policy query at source time $q\leq u$. Repeated
queries create a triangular cache of predictions. At physical time $t$, define
the fresh candidate
\begin{equation}
    F_t = E_{t,t}
\end{equation}
and, for a fixed delay $d=20$, the old candidate
\begin{equation}
    O_t = E_{t,t-20}.
\end{equation}
Both candidates target the same controller time; they differ only in the
observation time that conditioned their source query. With a 20-Hz controller,
$d=20$ is a 1.0-s source-observation gap.

We partition each action as $a=[a^{a},a^{g}]$, where $a^{a}\in\mathbb{R}^{6}$
contains relative translation and rotation commands and $a^{g}\in\mathbb{R}$
is the gripper command. The four factorial assignments are
\begin{align}
    \mathrm{FF}_t &= [F_t^{a},F_t^{g}], &
    \mathrm{OO}_t &= [O_t^{a},O_t^{g}], \\
    \mathrm{FO}_t &= [F_t^{a},O_t^{g}], &
    \mathrm{OF}_t &= [O_t^{a},F_t^{g}].
\end{align}
Before $O_t$ is available, all conditions execute $F_t$. No candidate is
averaged or smoothed in this factorial.

\subsection{Fixed asymmetric temporal reuse}

The confirmed executor is the fixed FO20 assignment
\begin{equation}
    a_t^{\mathrm{FO20}}=[F_t[0{:}6],O_t[6]], \qquad t\geq20.
    \label{eq:fo20}
\end{equation}
It has no learned parameters and does not modify the policy. Every step still
uses the current observation to issue one complete policy query and cache its
$100\!\times\!7$ prediction. FO20 changes only which cached source supplies
each component at execution. In particular, $O_t[6]$ advances from chunk
offset 20 of query $t-20$ to offset 20 of query $t-19$ on the next tick.
Therefore the executed gripper value may change every controller step.

The confirmatory baselines use the same query cache:
\begin{itemize}
    \item \textbf{Newest}: execute $F_t$ in full.
    \item \textbf{Full old20}: execute $O_t$ in full once available.
    \item \textbf{Age exponential}: average every valid full-action prediction
    for $t$ with shared weights $w_q\propto\exp[-0.03(t-q)]$ across all seven
    dimensions.
    \item \textbf{CogACT}: use the released cosine-similarity weighting with
    frozen $\alpha=0.3$ and one shared weight vector for all seven dimensions
    ~\cite{li2024cogact}.
\end{itemize}
The two weighted baselines test whether a strong shared full-action temporal
rule explains the benefit. The study does not search over $d$, $\beta$, or
$\alpha$.

\subsection{What a separable offline loss cannot measure}

Consider any teacher-forced loss that separates by component,
\begin{equation}
    L(a^{a},a^{g})=L_a(a^{a})+L_g(a^{g}).
\end{equation}
For a fixed target, write $A_F=L_a(F_t^a)$, $A_O=L_a(O_t^a)$,
$G_F=L_g(F_t^g)$, and $G_O=L_g(O_t^g)$. The four losses are
\begin{align}
L(FF)&=A_F+G_F, & L(OO)&=A_O+G_O,\\
L(FO)&=A_F+G_O, & L(OF)&=A_O+G_F.
\end{align}
Their balanced interaction is exactly
\begin{equation}
L(FF)+L(OO)-L(FO)-L(OF)=0.
\label{eq:separable_identity}
\end{equation}
This algebraic identity does not imply that the four actions are behaviorally
equivalent. It states only that an additive teacher-forced loss contains no
arm--gripper interaction term with which to score symmetric composition.
